\section{Introduction}
\label{sec:intro}

Steering vectors have become a standard tool for controlling language model behavior at inference time \citep{turner2023steering, rimsky2024steering}. The core idea is simple: identify a direction in activation space that corresponds to a concept, then add a scaled version of that direction during generation to shift output toward the target concept. This approach rests on a strong assumption---that concepts are encoded as linear directions in the model's residual stream. The Linear Representation Hypothesis (\lrh) formalizes this view, and substantial evidence supports it for many concept types \citep{park2023linear, jiang2024origins}.

Yet a persistent puzzle remains. Some styles and behaviors resist linear steering despite appearing to be ``known'' by the model. \citet{konen2024style} find that disgust achieves near-random steering accuracy (AUC~0.51), and surprise yields few successful steering vectors, even though language models can readily produce disgusted or surprised text when prompted. If the model can generate a style on demand via prompting, why would a linear intervention in activation space fail? One possibility is that these styles are encoded \emph{non-linearly}, making them invisible to linear probes and steering vectors alike. Another is that the representation is linear but the \emph{causal pathway} from representation to generation is not.

We disentangle these possibilities by conducting the first systematic comparison of three facets of style encoding across a broad taxonomy of 11 emotional styles: (i)~\emph{linear separability}, measured by comparing linear and MLP probes; (ii)~\emph{steerability}, measured by contrastive activation addition (\caa); and (iii)~\emph{promptability}, measured by style adherence ratings. We test these on \pythia using the \goemotions dataset \citep{demszky2020goemotions}.

Our central finding is a \emph{detection--control gap}: all 11 styles are linearly separable in activation space (non-linearity indices $\approx 1.0$), yet steering effectiveness varies dramatically and is uncorrelated with linear separability ($r = -0.220$, $p = 0.52$). Meanwhile, prompting succeeds uniformly for all styles (100\% success rate, 3.8--4.4/5 adherence), including those that steer poorly. The correlation between steering and prompting success is essentially zero ($r = -0.018$, $p = 0.96$).

These results have direct implications for AI safety and controllable generation. Linear intervention methods---steering vectors, sparse autoencoder features---may systematically miss behaviors that are linearly detectable but causally accessed through nonlinear pathways. Our contributions are:
\begin{itemize}[leftmargin=*,itemsep=0pt,topsep=0pt]
    \item We show that all 11 tested emotional styles are linearly represented in \pythia (MLP probes never outperform linear probes), supporting the \lrh for representation geometry.
    \item We demonstrate that linear separability does not predict steering effectiveness, revealing a detection--control gap where representation geometry and causal intervention diverge.
    \item We establish that prompting and steering access fundamentally different pathways to style control, with prompting succeeding uniformly where steering fails selectively.
\end{itemize}

The rest of this paper is organized as follows. \Secref{sec:related} reviews related work on steering vectors, the linear representation hypothesis, and non-linear representations. \Secref{sec:method} describes our experimental methodology. \Secref{sec:results} presents our main results. \Secref{sec:discussion} discusses implications and limitations. \Secref{sec:conclusion} concludes.
